\documentclass[11pt]{article}
\usepackage[margin=25mm]{geometry}
\usepackage{fontspec,unicode-math}
\setmainfont{Latin Modern Roman}
\setmathfont{Latin Modern Math}
\usepackage{amsmath,amssymb,mathtools}
\usepackage{xurl,hyperref}
\hypersetup{colorlinks=true,urlcolor=blue,linkcolor=blue}
\setcounter{secnumdepth}{0}
\setlength{\parindent}{0pt}
\setlength{\parskip}{5pt}
\setlength{\emergencystretch}{3em}
\newcommand{\sref}[2]{\hyperlink{source:#1:#2}{[S#2]}}
\newcommand{\cataloguescope}[1]{\paragraph{Source question.}#1}
\newcommand{\floor}[1]{\left\lfloor#1\right\rfloor}
\begin{document}
\section{Complexity of Bray’s Involution-Centralizer Algorithm}\label{statement}

\subsection{Definitions and mathematical statement}

\paragraph{Shared notation.}Standard number systems are denoted by $\mathbb N,\mathbb Z,\mathbb Q,\mathbb R,\mathbb C$; index ranges, computational models and specialized notation follow the cited problem source and the definitions below.


A strong involution $t$ has order two and both its $+1$ and $-1$ eigenspaces have dimensions between $n/3$ and $2n/3$. Bray\textquotesingle s method samples independent uniform $g\in G$, sets $y=tt^g$ with $t^g=g^{-1}tg$, and extracts centralizer elements from the odd/even order cases. The centralizer is $C_G(t)=\{h:ht=th\}$.

The group $\mathrm{GL}_n(\mathbb Z)$ consists of integer matrices with integer inverses, and $\mathrm{SL}_n$ imposes determinant one. The subgroup generated by matrices consists of finite products of them and their inverses. A centralizer consists of elements commuting with the specified element.
\paragraph{Statement recorded in the supplied source.}
For a finite classical group of dimension $n$ over a field of odd characteristic, does Bray's algorithm for constructing the centralizer of a strong involution succeed with high probability after examining only $O(\log n)$ random elements? \sref{statement}{2}
\paragraph{Scope qualification.}
The original report treats the even-case construction of a subgroup containing the second derived subgroup $C_G(t)^{\prime\prime}$ as sufficient for its modified aim. This is not silently strengthened to generating every element of the full centralizer.
\paragraph{Source status.}
The supplied catalogue labels this record partially solved. Its source-status metadata was checked against the supplied snapshot on 2026-10-11. This import does not independently establish current open status or mathematical correctness. The catalogue literature assessment is dated 2026-08-21.
\subsection{Short English statement}
Prove logarithmic random-sampling complexity for the stated strong-involution centralizer algorithm.
\subsection{Sources}
\begin{itemize}
\item[\hypertarget{source:statement:1}{S1}] Ulam AI and the UnsolvedMath Contributors, UnsolvedMath: A Curated Collection of Open Mathematics Problems, catalogue entry. Catalogue curation and metadata: CC BY 4.0. The supplied source question is retained; definitions and the English summary are editorial additions.
\url{https://huggingface.co/datasets/ulamai/UnsolvedMath}.
\item[\hypertarget{source:statement:2}{S2}] Computational Group Theory (2017), Oberwolfach Reports, DOI 10.4171/owr/2016/37, pp. 2123--2169
\url{https://doi.org/10.4171/owr/2016/37}.
\item[\hypertarget{source:statement:3}{S3}] Dixon--Praeger--Seress
\url{https://www.sciencedirect.com/science/article/pii/S002186931730635X}.
\item[\hypertarget{source:statement:4}{S4}] Glasby--Praeger--Roney-Dougal
\url{https://doi.org/10.1016/j.jalgebra.2019.09.009}.
\item[\hypertarget{source:statement:5}{S5}] Original Oberwolfach Report 37/2016, Cheryl E. Praeger, pp. 2160--2161
\url{https://ems.press/content/serial-article-files/46642}.
\textit{Underlying Oberwolfach Reports material retains its source-specific licence, including CC BY-SA 4.0 where applicable; see THIRD\_PARTY\_NOTICES.md.}
\end{itemize}
\paragraph{Preparation and review.}
Original question attributed to the cited source and UnsolvedMath contributors. Codex (GPT-6; exact serving revision not exposed) assisted with definitions, English summaries, source inspection and formatting. No new human mathematical authorship or independent proof verification is asserted. Catalogue identity checked against the complete supplied registry and existing repository statements; known flagged repetitions and obvious identical conjectures excluded. Full mathematical-equivalence/literature review has not been performed. Statement preparation and final review followed contribution prompts C and D; automated checks verify structure and identity, not open status.
\label{end:statement}
\end{document}
